\begin{center}
\large \textsc{Biographical sketch} --- Kaloyan M.~Penev
\end{center}
%
\vspace{-3mm}
% +++++++++++++++++++++++++++++++++++++++++++++++++++++++++++++++++++++
\begin{comment}

NASA:
http://www.hq.nasa.gov/office/procurement/nraguidebook/
The PI (and Co-PI) must include a biographical sketch (not to exceed two
pages) that includes his/her professional experiences and positions and a
bibliography of recent publications, especially those relevant to the
proposed investigation. A one-page sketch for each Co-Investigator must also
be included (Note: Any Co-I also serving in one of the three special Co-I
categories defined in Section 1.4.2 may use the same two-page limit as for
the PI). For the PI and any Co-Is who are required to provide Current and
Pending Support information (ref. Section 2.3.8), the biographical sketch
must include a description of scientific, technical and management
performance on relevant prior research efforts. Those participants who will
play critical management or technical roles in the proposed investigation
should demonstrate that their qualifications, capabilities, and experience
are appropriate to provide confidence that the proposed objectives will be
achieved

NSF:
http://www.nsf.gov/pubs/policydocs/pappguide/nsf13001/gpg_2.jsp#IIC2f

\end{comment}
% +++++++++++++++++++++++++++++++++++++++++++++++++++++++++++++++++++++

% #####################################################################
% #####################################################################
%
{\bf \noindent Professional Preparation\\}
\vspace{-3mm}

\begin{tabular}{p{60mm}p{70mm}p{20mm}}
	\hline\hline
	{\em Undergraduate Institution}	&	Major	&	Degree\\
	California Institute of Technology	&
	Astronomy	&
	BSc, 2003\\
%
	{\em Graduate Institution}	&	&\\
	Harvard University	&
	Astronomy	&
	PhD, 2009\\
%
	{\em Postdoctoral Institution(s)}	&	Area & Incl.~dates\\
	Harvard Smithsonian Center for Astrophysics &
	HATSouth \& HATNet surveys	&
	2009\---2011\\
\hline

\end{tabular}

% #####################################################################
% #####################################################################
%
\vspace{2mm}
%
{\bf \noindent Appointments\\}
%
\vspace{-5mm}
%
\begin{itemize}
%
\item 2017 --- present: {\em Assistant Prof.}, University of Texas at Dallas,
Department of Physics.
%
\item 2011 --- 2017: {\em Associate Research Scholar}, Princeton University,
Department of Astrophysics.
%
\item 2009 --- 2011: {\em Postdoctoral Fellow}, Harvard-Smithsonian
Center for Astrophysics (CfA).
\end{itemize}

\vspace{1mm}
% #####################################################################
% #####################################################################
%
{\bf \noindent Performance on relevant prior research efforts:\\}
%
\begin{itemize}
%
        \vspace{-5mm}
%
    \item K.~Penev is the PI on NASA ATP grant 80NSSC18K1009 (2018--2022),
        ``Comprehensive Empirical Constraints on Tidal Dissipation Rates : Q is
        Not a Constant''. So far, work under this grant has resulted in four
        published articles, four conference proceedings, and on-going work on 2
        additional articles fully funded by this grant. In addition with support
        from this grant K.P. contributed to 9 additional publications. The work
        plan of the proposal called for 5 publications and two conference
        presentations.
%
    \item K.~Penev was the PI on NASA Origins grant NNX13AQ62G (2013-2015),
        ``The Origins of Close-in Extrasolar Planets''. Exclusively under this
        effort, K.P. developed the P.O.E.T. code, lead two papers, and two
        conference presentations, advised two undergraduate students, one of
        which lead one of the publications, and is a co-author on another.
        Further, K.P. contributed to 22 additional articles. K.P.  advised two
        undergraduate students, one of which lead one of the publications, and
        is a co-author on another. Three papers were expected when proposal was
        submitted and no student involvement.
%
    \item K.~Penev was funded under NASA Origins grant NNX09AB29G (2009-2013),
        ``Operating a Global Network of Southern Hemisphere Automated Telescopes
        to Detect and Characterize a Multitude of Transiting Exoplanets''. K.P.
        is the author of the calibration pipeline that processes all the data of
        the HAT-South and HATNet networks as well as a number of developments
        which significantly improved performance. As part of this effort, K.P.
        is an author in 9 publications based on HATSouth data, 9 based on HATNet
        data.
%
\end{itemize}

% #####################################################################
% #####################################################################
%
\vspace{1mm}
%
% #####################################################################
% #####################################################################
%
\vspace{1mm}
{\bf \noindent Relevant Scientific Publications:\\}
\vspace{-5mm}
\begin{enumerate}
%

%%%% %T\n %%%% %u\n \\item %\l~%\Y, \n {``\\em %\T''}, {\\it %\j}, {\\bf %\V}, %\p\n,  %\X
%%%% %T\n %%%% %u\n \\item %\8.8l~%\Y, \n {``\\em %\T''}, {\\it %\j}, {\\bf %\V}, %\p\n,  %\X

% ++++++++++++++++++
    \item \textbf{Penev, K.}; Schussler, J. ``Comprehensive Bayesian modelling
        of tidal circularization in open cluster binaries part I: M 35, NGC
        6819, and NGC 188'' 2022, {\em MNRAS} 516,6145

    \item \textbf{Penev, K.}; Bouma, L.~G.; Winn, J.~N.; Hartman, J.~D.
        ``Empirical Tidal Dissipation in Exoplanet Hosts From Tidal Spin-up''
        2018, {\em AJ} 155, 165

    \item \textbf{Penev, K.}; Hartman, J.~D.; Bakos, G.~{\'A}; and 22
        coauthors ``HATS-18b: An Extreme Short-period Massive Transiting
        Planet Spinning Up Its Star'' 2016, {\em AJ} 152, 127

	\item \textbf{Penev, K.}; Zhang, M; Jackson, B. ``POET: A Model for
		(P)lanetary (O)rbital (E)volution due to (T)ides on Evolving Stars''
        2014, {\em PASP} 126, 553

	\item \textbf{Penev, K.}; Bakos, G.~\'A.; Bayliss, D.; Jordan, A.;
		Mohler, M.; and 23 coauthors
		``HATS-1b: The First Transiting Planet Discovered by the HATSouth
		Survey'' 2013, {\em AJ} 145, 5

	\item \textbf{Penev, K.}; Jackson, B., Spada, F., Thom, N.
		``Constraining Tidal Dissipation in Stars from The Destruction Rates
		of Exoplanets'' 2012, {\em ApJ} 751, 96

	\item \textbf{Penev, K.}; Barranco, J.; Sasselov, D. ``Three-dimensional
		Spectral Simulations of Anelastic Turbulent Convection'' 2011, {\em
		ApJ} 734, 118

	\item \textbf{Penev, K.}; Sasselov, D. ``Tidal Evolution of Close-in
		Extrasolar Planets: High Stellar Q from New Theoretical Models''
		2011, {\em ApJ} 731, 67

	\item \textbf{Penev, K.}; Barranco, J.; Sasselov, D. ``Direct Calculation
		of the Turbulent Dissipation Efficiency in Anelastic Convection''
		2009, {\em ApJ} 705, 285

	\item \textbf{Penev, K.}; Sasselov, D.; Robinson, F.; Demarque, P.
		``Dissipation Efficiency in Turbulent Convective Zones in Low-Mass
		Stars'' 2009, {\em ApJ} 704, 930

	\item \textbf{Penev, K.}; Sasselov, D.; Robinson, F.; Demarque, P. ``On
		Dissipation inside Turbulent Convection Zones from Three-dimensional
		Simulations of Solar Convectio'' 2007, {\em ApJ} 655, 1166

    \item Patel, R; \textbf{Penev, K.} ``Constraining tidal quality factor using
        spin period in eclipsing binaries'' 2022, {\em MNRAS} 512, 3651

    \item Anderson, K; Winn, J; \textbf{Penev, K.} ``On a Possible Solution to
        the Tidal Realignment Problem for Hot Jupiters'' 2021, {\em ApJ} 914, 56

    \item Jordan, A.; Hartman, J. D.; Bayliss, D.; and 15 co-authors
        ``HATS-74Ab, HATS-75b, HATS-76b, and HATS-77b: Four Transiting Giant
        Planets Around K and M Dwarfs'' 2022, {\em AJ} 163, 125

    \item Sha, L.; Huang, C.; Shporer, A. and 70 co-authors ``TOI-954 b and
        K2-329 b: Short-period Saturn-mass Planets that Test whether Irradiation
        Leads to Inflation'' 2021, {\em AJ} 161, 82

    \item Bakos, G.; Hartman, J.; Bhatti, W.; Csubry, Z.; \textbf{Penev, K.} and
        29 co-authors ``HAT-P-58 - b-HAT-P-64b: Seven Planets Transiting Bright
        Stars'' 2021, {\em AJ} 162, 7

    \item Lindor, B.; Hartman, J.; Bakos, G.; Bhatti, W.; Csubry, Z.;
        \textbf{Penev, K.} and 20 co-authors ``HAT-P-68b: A Transiting Hot
        Jupiter around a K5 Dwarf Star'' 2021, {\em AJ} 161, 64

    \item Jordan, A.; Bakos, G.; Bayliss, D.; and 24 co-authors ``HATS-37Ab and
        HATS-38b: Two Transiting Hot Neptunes in the Desert'' 2020, {\em AJ},
        160, 222

    \item Rodriguez M.; Gaudi, S.; Rodriguez, J.; Zhou, G.; Labadie-Bartz, J.;
        Quinn, S.; \textbf{Penev, K.} and 85 co-authors ``KELT-25 b and KELT-26
        b: A Hot Jupiter and a Substellar Companion Transiting Young A Stars
        Observed by TESS'' {\em AJ} 160, 111

    \item Bakos, G.; Hartman, J.; Bhatti, W.; Csubry, Z.; \textbf{Penev, K.} and
        29 co-authors ``HAT-P-58b -- HAT-P-64b: Seven Planets Transiting Bright
        Stars'' {\em arXiv:2007.05528}

    \item Bakos, G.; Bayliss, D.; Bento, J. and 41 co-authors ``HATS-71b: A
        Giant Planet Transiting an M3 Dwarf Star in TESS Sector 1'' {\em AJ}
        159, 267

    \item Hartman, J. ; Jordán, A.; Bayliss, D. and 43 co-authors ``HATS-47b,
        HATS-48Ab, HATS-49b, and HATS-72b: Four Warm Giant Planets Transiting K
        Dwarfs'' 2020, {\em AJ} 159, 173

    \item Mancini, L.; Sarkis, P.; Henning, Th.; Bakos, G. Á.; Bayliss, D.;
        Bento, J.; Bhatti, W.; Brahm, R.; Csubry, Z.; Espinoza, N.; Hartman, J.;
        Jordán, A.; \textbf{Penev, K.}; and 8 co-authors ``The highly inflated
        giant planet WASP-174b'' 2020, \em{A\&A} 633, 30

    \item Rodriguez, J.; Eastman, J.; Zhou, G.; Quinn, S.; Beatty, T.;
        \textbf{Penev, K.}; and 73 co-authors ``KELT-24b: A 5M J Planet on a 5.6
        day Well-aligned Orbit around the Young V = 8.3 F-star HD 93148'' 2019,
        \em{AJ} 158, 197

    %\item Zhou, G.; Huang, C.; Bakos, G.; Hartman, J.; Latham, D.; Quinn, S.;
    %    Collins, K.; Winn, J.; Wong, I.; Kovács, G.; Csubry, Z.; Bhatti, W.;
    %    \textbf{Penev, K.}; and 58 co-authors ``Two New HATNet Hot Jupiters
    %    around A Stars and the First Glimpse at the Occurrence Rate of Hot
    %    Jupiters from TESS'' 2019, \em{AJ} 158, 141

\end{enumerate}
